\paragraph{Contemporaneous factorial control.}
\label{sec:eval:factorial}
Four generation settings cross preprocessing ($P$) and diagnostic context
($D$) on the same 60 controlled and 60 natural inputs (480 Gemma responses).
Each response is scored with and without the same gate ($G$). Main effects
average over the other two switches, separating component contributions
from method-specific prompt differences.

\begin{table}[t]

\centering

\caption{Factorial main effects in F1 percentage points, averaged over the other two switches. Each cohort has sixty paired documents. Intervals are marginal 95\% document-bootstrap intervals; $p_H$ uses Holm correction over six main effects.}

\label{tab:recovery_factorial}

\small

\begin{tabular}{llrrr}

\toprule

Input & Component & Change & 95\% interval & $p_H$ \\

\midrule

Controlled & Preprocessing & -0.028 & [-0.797, 0.778] & 1.000 \\

Controlled & Diagnosis & -0.651 & [-1.348, -0.056] & 0.164 \\

Controlled & Gate & +0.164 & [0.054, 0.301] & 0.156 \\

Natural & Preprocessing & +0.697 & [0.000, 1.905] & 1.000 \\

Natural & Diagnosis & -0.535 & [-2.376, 0.842] & 1.000 \\

Natural & Gate & +0.707 & [0.225, 1.316] & 0.086 \\

\bottomrule

\end{tabular}

\end{table}

None of the six main effects passes Holm correction (Table~\ref{tab:recovery_factorial}). Diagnostic context has a negative mean effect in both cohorts. Gate improves mean F1 by 0.164 points on controlled inputs and 0.707 points on natural inputs. Its full-check audit rejects 47 candidate occurrences across all cohorts: 17 have wrong heads and 30 lack source support; none exactly matches a reference triple. Interactions and per-check removals are released with the paired scores.

These small filtering effects and the six nonsignificant adjusted tests
support a limited candidate-filtering role; they do not establish a stable
quality gain from every component.
